\section{Comparison and cancellation}\label{sec:transfer}

We now establish the two estimates that will control the quantity
$Y_l(N,A)$ in the induction.  The first compares a lift, after some prime
coordinates have been fixed, with lifts of sets on shorter intervals.  The
second gives a power saving when the small-prime residues permit every
increment of the designated cycle.  The progression comparison and the
reduction to an ordered pair adapt the arguments of
\cite[Sections~7--8, especially Propositions~7.2 and~8.1]{squarepower}.
For a modulus $D$, our progression spacing is $\lambda(D)$, and the
change of variables also changes the auxiliary parameter from $l$ to $lD$.

For positive integers $l,N$, let
\[
 \mathcal Y_l(N)=\max_A Y_l(N,A),
\]
where the maximum is over all $A\subseteq[N]$ satisfying
\begin{equation}\label{eq:24}
 (A-A)\cap h_l(\N_+)\subseteq\{0\}.
\end{equation}
We retain the parameters of \eqref{eq:14}, in particular
$Q=N^\beta$, $H=N^{\beta/2}$, and
$\mathcal P_N=\{p\text{ prime}:P\le p\le Q\}$.
For a fixed set $A\subseteq[N]$, define
\[
 f_s(x)=M_0\1_{\{x\equiv s\pmod{M_0}\}}\1_A(x),
 \qquad
 G_s=\mathcal L_{[N],\mathcal P_N,Q}f_s,
 \qquad s\in\Z/M_0\Z.
\]
Thus $\1_A$ is the uniform average of the functions $f_s$, and its lift
is the corresponding average of the $G_s$.

\subsection{Fourier coefficients after specialization}

Fixing prime coordinates restricts the underlying integers to a residue
class, provided the cutoff includes all characters on those coordinates.
The precise identity is the following.

\begin{lemma}[Specialization of a lift]\label{transfer:specialization}
Let $I\subseteq[N]$ be a nonempty integer interval, let
$B\subseteq\mathcal P_N$, and put $D=M_0q_B$.
Fix $s\in\Z/M_0\Z$ and $b=(b_p)_{p\in B}\in X_B$.
Let $g$ be the function on $X_{\mathcal P_N\setminus B}$ obtained by
specializing $\mathcal L_{I,\mathcal P_N,Q}f_s$ at $z_B=b$.
For every remaining-coordinate frequency $\xi$ with
$q(\xi)q_B\le Q$, one has
\begin{equation}\label{eq:25}
 \widehat g(\xi)
 =\frac{D}{|I|}
   \sum_{\substack{x\in I\\x\equiv a'\pmod D}}
       \1_A(x)e(-x\xi),
\end{equation}
where $a'\pmod D$ is determined by
$a'\equiv s\pmod{M_0}$ and $a'\equiv b_p\pmod p$ for $p\in B$.
\end{lemma}

\begin{proof}
The frequencies contributing to $\widehat g(\xi)$ have the form
$\xi+\omega$, where $\omega$ ranges over the characters on $X_B$
that survive the cutoff.  The denominators are coprime, so
$q(\xi+\omega)=q(\xi)q(\omega)\le q(\xi)q_B\le Q$.
Consequently every character on $X_B$ occurs in the specialized
coefficient, which equals
\[
 \frac{M_0}{|I|}\sum_{x\in I}
 \1_A(x)\1_{\{x\equiv s\pmod{M_0}\}}e(-x\xi)
 \sum_{\omega\text{ on }X_B}e((b-x)\omega).
\]
Character orthogonality makes the inner sum $q_B$ when
$x\equiv b_p\pmod p$ for every $p\in B$, and zero otherwise.
The Chinese remainder theorem now gives \eqref{eq:25}.
\end{proof}

\subsection{Comparison with a shorter interval}

The next estimate is uniform in the auxiliary parameter $l$.  This
uniformity is necessary because a progression replaces $l$ by $lD$, which
need not satisfy the size restriction used in the kernel estimate.

\begin{proposition}[Shorter-scale comparison]\label{transfer:shorter-scale}
Let $N$ be sufficiently large, let $l\ge1$, and let $A\subseteq[N]$
satisfy \eqref{eq:24}.  Suppose that $B\subseteq\mathcal P_N$ and
$q_B\le N^\tau$.  Set
\begin{gather*}
 \mathcal S=\mathcal P_N\setminus B,
 \qquad D=M_0q_B,
 \qquad L_D=\lambda(D),\\
 n_B=\left\lceil\frac{N}{L_D}\right\rceil,
 \qquad \Lambda_B=\frac{L_Dn_B}{N}.
\end{gather*}
For each $v\in V$, choose $s_v\in\Z/M_0\Z$ and specialize $G_{s_v}$
at an arbitrary vector of coordinates on $B$, obtaining a function $g_v$
on $X_{\mathcal S}$.  Then
\begin{equation}\label{eq:26}
 \bigl|\mathcal Z_{l,\mathcal S}((g_v)_{v\in V})\bigr|
 \le \Lambda_B^d\mathcal Y_{lD}(n_B)+O(N^{-\sigma}).
\end{equation}
The scale and the rounding factor satisfy $1\le n_B<N$ and
\begin{equation}\label{eq:27}
 1\le\Lambda_B\le1+M_0^kN^{k\tau-1}.
\end{equation}
The implied constant and the threshold for $N$ are independent of
$l,A,B$ and all specializations.
\end{proposition}

\begin{proof}
We first record the scale inequalities.  By \eqref{eq:1},
$D\le L_D\le D^k$, so $L_D\ge M_0\ge2$ and $n_B<N$ for large $N$.
The defining ceiling gives \eqref{eq:27}.  Also
\begin{gather*}
 n_B\ge M_0^{-k}N^{1-k\tau}\ge N^{1/2},\\
 H\le Q':=n_B^\beta\le Q,
 \qquad q_BH\le Q.
\end{gather*}
Here $k\tau<1/2$ and $\tau<\beta/2$ follow from \eqref{eq:14}.

By the signed H\"older inequality \eqref{eq:10}, it is enough to prove
the asserted common bound for $Z_{l,\mathcal S}(g)$, where $g$ is any
one of the specialized inputs.  Let $g_{\le H}$ retain exactly the
remaining-coordinate frequencies of denominator at most $H$.
We claim that
\[
 Z_{l,\mathcal S}(g)
 =Z_{l,\mathcal S}(g_{\le H})
   +O(H^{-\gamma}q_BN^{o(1)}).
\]
To see this, telescope the $d$ slots and apply \eqref{eq:11} to the
high-frequency input in each term.  The moment estimate \eqref{eq:12}
is applicable to all the other inputs as well as to that high-frequency
piece.  Indeed, before specialization, a cut on the remaining denominator
retains entire supports $C\subseteq\mathcal P_N$ according to the
condition on $q_{C\setminus B}$; it never selects only part of the
frequencies with a given full support $C$.  Fixing $B$ costs at most
$q_B^{1/2}$ in the distinguished $L^2$ norm and
$q_B^{1/[2(d-1)]}$ in each of the other $d-1$ norms.  The total cost is
therefore $q_B$.  The original input is bounded by $M_0$, and
$N\ge10Q^6$ for large $N$, so \eqref{eq:12} supplies the remaining
$N^{o(1)}$ factor.  Since
\[
 \frac{\beta\gamma}{2}-\tau-\sigma
 =\tau\left(500k(d+1)-1-\frac\gamma4\right)>0,
\]
the claimed error is $O(N^{-\sigma})$.  The denominator cuts are
symmetric under negation, so the functions used here are real.

We next express $g_{\le H}$ as an average of shorter lifts.  Let
$a'\pmod D$ be the residue in Lemma~\ref{transfer:specialization}.
For each representative $j\in[L_D]$ with $j\equiv a'\pmod D$, set
\[
 c_j=j-L_D,\qquad
 A'_j=\{x'\in[n_B]:L_Dx'+c_j\in A\}.
\]
These progressions partition the integers of $[N]$ in the class
$a'\pmod D$: the first positive term is $j$, and $n_B$ terms suffice
for every such integer.  Terms above $N$ contribute no points to $A'_j$.
Moreover, $A'_j$ satisfies \eqref{eq:24} with $lD$ in place of $l$.
For if a nonzero difference in $A'_j-A'_j$ were $h_{lD}(m)$ with
$m\ge1$, the corresponding difference in $A-A$ would be
\[
 L_Dh_{lD}(m)=h_l(Dm-i),\qquad 0\le i<D,
\]
by \eqref{eq:1}; its argument is positive and its value is nonzero.

Put $\mathcal H_j=\mathcal L_{[n_B],\mathcal S,H}\1_{A'_j}$.
Every prime of $L_D$ divides $D$, so $L_D$ is invertible on
$X_{\mathcal S}$ and multiplication by $L_D$ preserves denominators.
Lemma~\ref{transfer:specialization}, applied using $q_BH\le Q$, gives
the exact identity
\begin{equation}\label{eq:28}
 g_{\le H}(z)
 =\frac{\Lambda_B}{L_D/D}
   \sum_{\substack{j\in[L_D]\\j\equiv a'\pmod D}}
    \mathcal H_j\bigl(L_D^{-1}(z-c_j)\bigr).
\end{equation}
For completeness, the summand on the right has Fourier coefficient
$\widehat{\mathcal H_j}(L_D\xi)e(-c_j\xi)$ at $\xi$.
Since $\Lambda_B/(L_D/D)=Dn_B/N$, summing these coefficients gives
$D/N$ times the sum over the required original integers, exactly as in
\eqref{eq:25}; both sides have zero coefficients above $H$.

Under the map $z\mapsto L_D^{-1}(z-c_j)$, the tuple law for $l$ pushes
forward to the tuple law for $lD$, by \eqref{eq:7} at each remaining
prime.  There are exactly $L_D/D$ representatives $j$.
The seminorm triangle inequality applied to \eqref{eq:28}
therefore yields
\[
 Z_{l,\mathcal S}(g_{\le H})
 \le\Lambda_B^d\max_j Z_{lD,\mathcal S}(\mathcal H_j).
\]
This is where the exact change of tuple laws preserves the leading
constant in the comparison.

It remains to replace the cutoff $H$ on the shorter interval by its
natural cutoff $Q'$.  Equations~\eqref{eq:11} and~\eqref{eq:12}, now
on $[n_B]$, bound the change of the diagonal by
$O(H^{-\gamma}N^{o(1)})=O(N^{-\sigma})$.
Their interval hypothesis holds because $6\beta<1$, and hence
$n_B\ge10(Q')^6$ for large $N$.
The extended lift only uses primes in $\mathcal P_{n_B}$.
It is obtained from
$\mathcal L_{[n_B],\mathcal P_{n_B},Q'}\1_{A'_j}$ by averaging the
coordinates in $B\cap\mathcal P_{n_B}$ and then viewing the result
as independent of the additional coordinates in $\mathcal S$.
The contraction under coordinate averaging from
Lemma~\ref{lift:seminorm} gives
\[
 Z_{lD,\mathcal S}(\mathcal H_j)
 \le Y_{lD}(n_B,A'_j)+O(N^{-\sigma})
 \le\mathcal Y_{lD}(n_B)+O(N^{-\sigma}).
\]
Combining these estimates with \eqref{eq:27} proves the diagonal bound.
Applying \eqref{eq:10} to the mixed inputs proves \eqref{eq:26}:
the product of $d$th roots of quantities bounded by the same nonnegative
upper bound is at most that upper bound.
\end{proof}

\subsection{Cancellation for good residue assignments}

The previous comparison applies to every residue assignment.  We now
use avoidance to obtain a stronger estimate for assignments compatible
with the designated cycle.  An assignment
$(s_v)_{v\in V}\in(\Z/M_0\Z)^V$ is \emph{good for $l$} if every
increment $s_{v_{j+1}}-s_{v_j}$, $0\le j<L$, is admissible for $l$
modulo $M_0$.  As before, $v_L=v_0$.

\begin{lemma}\label{transfer:good-assignments}
For every $l\ge1$, the proportion $\rho_l$ of good assignments under
independent uniform choices of all the $s_v$ satisfies
\[
 \rho_l\ge\rho_*:=M_0^{1-L}>0.
\]
\end{lemma}

\begin{proof}
At every prime-power factor of $M_0$, the construction in
Section~\ref{sec:local} supplies a length-$L$ sequence of admissible
residues whose sum is zero.  The Chinese remainder theorem combines them
into such a sequence modulo $M_0$.
Fix one sequence.  There are $M_0$ choices of the initial cycle residue;
the other cycle residues are then determined, and the zero-sum condition
closes the cycle.  Its $L$ vertices are distinct.  The remaining $d-L$
residues can be arbitrary, giving at least $M_0^{d-L+1}$ good assignments
out of $M_0^d$.
\end{proof}

The proportion $\rho_l$ may vary with $h$ and $l$. When
Condition~\eqref{local:unit-condition} holds and the designated choices of
Sections~\ref{sec:local} and~\ref{sec:tuples} are made, both $M_0$ and $L$,
and therefore the lower bound $\rho_*=M_0^{1-L}$, depend only on $k$.

\begin{proposition}[Cancellation on good assignments]\label{transfer:cancellation}
Let $N$ be sufficiently large, let $l\le N^\zeta$, and let
$A\subseteq[N]$ satisfy \eqref{eq:24}.  If $(s_v)_{v\in V}$ is good
for $l$, then
\begin{equation}\label{eq:29}
 \left|
  \int\prod_{v\in V}G_{s_v}(z_v)\,
    d\bigotimes_{p\in\mathcal P_N}
       (\mu_{l,p}-\eps_{l,p})
 \right|\ll N^{-\sigma}.
\end{equation}
The bound is uniform in $l,A$ and the good assignment.
\end{proposition}

\begin{proof}
We may replace each $\mu_{l,p}-\eps_{l,p}$ by its normalized law
$\mu_{l,p}^{\circ}$: passing back multiplies the integral by
$\prod_{p\in\mathcal P_N}(1-\eta_{l,p})\le1$.
We first locate an ordered pair of intervals along the cycle.
We then truncate the other inputs and condition on their prime supports;
the remaining pair integral will have exactly the multipliers appearing
in Proposition~\ref{kernel:estimate}.

Partition $[N]$ into $\lfloor N^\tau\rfloor$ consecutive intervals
whose lengths differ by at most one.  Each $G_s$ is the average of the
interval lifts $\mathcal L_{I,\mathcal P_N,Q}f_s$ with weights $|I|/N$.
Thus we may average independently over an interval assignment to every
slot.  These intervals have length $\gg N^{1-\tau}$, so the moment
estimate \eqref{eq:12} applies: indeed
$6\beta+\tau<7\beta<1$ by \eqref{eq:14}.
For any interval assignment, ordinary H\"older, uniform single-slot
marginals and \eqref{eq:12} bound the integral in absolute value by
$N^{o(1)}$.
The probability that all $L$ cycle slots choose the same interval is
\[
 \sum_I\left(\frac{|I|}{N}\right)^L
 \le\left(\max_I\frac{|I|}{N}\right)^{L-1}
 \ll N^{-(L-1)\tau}.
\]
Their contribution is $O(N^{-\sigma})$, since
$(L-1)\tau>\sigma=\gamma\tau/4$.

In every other assignment, some directed cycle edge has its first
interval $I_1$ strictly before its second interval $I_2$.
Otherwise the interval indices would be nonincreasing all the way
around the cycle, and therefore constant.  Fix such an edge for the
current assignment.  The constants below are uniform in this choice.

Set $T=N^\tau$ and truncate every interval lift except the two selected
ones to denominators at most $T$.
Telescoping the other slots, \eqref{eq:11} and \eqref{eq:12} give total
error $O(T^{-\gamma}N^{o(1)})=O(N^{-\sigma})$, because
$\gamma\tau>\sigma$.
Expand the truncated inputs into characters.  Each coefficient has
modulus at most $M_0$, and there are at most $T^2$ rational frequencies
of denominator at most $T$.  Thus the sum of the absolute coefficients
in the resulting expansion is $O(T^{2(d-2)})$.

Fix one term, let $B$ be the union of the supports of its $d-2$
characters, and put
\[
 R=q_B\le T^{d-2},\qquad
 \mathcal S=\mathcal P_N\setminus B.
\]
Condition on all tuple labels at every prime in $B$.
The expanded characters become constants of modulus one.
By independence between primes, the remaining tuple labels still have
law $\bigotimes_{p\in\mathcal S}\mu_{l,p}^{\circ}$.
Let $g_1,g_2$ denote the two selected interval lifts after specialization
at this conditioning.  They are functions on $X_{\mathcal S}$.
For the chosen directed edge, let $\Gamma$ be the increment multiplier
defined in Section~\ref{sec:kernel}.  The remaining pair integral is
exactly
\begin{equation}\label{eq:30}
 \sum_{\xi\text{ on }X_{\mathcal S}}
      \widehat g_1(-\xi)\widehat g_2(\xi)\Gamma(\xi).
\end{equation}
Indeed, simultaneous translation of the two labels kills all pairs of
characters whose frequencies do not sum to zero.  The surviving pair
with frequencies $-\xi,\xi$ has value
$e(\xi(z_2-z_1))$, whose expectation is $\Gamma(\xi)$.

For the terms with $q(\xi)>H$, \eqref{eq:15} and Parseval give
\[
 \left|\sum_{q(\xi)>H}
      \widehat g_1(-\xi)\widehat g_2(\xi)\Gamma(\xi)\right|
 \le H^{-\gamma}\|g_1\|_2\|g_2\|_2
 \ll H^{-\gamma}R N^{o(1)}.
\]
The last inequality uses the interval moment bound and the cost
$R^{1/2}$ of each specialization in $L^2$.

We apply the kernel estimate to the remaining terms.
Put $D=M_0R$.  The parameter choice gives
\[
 (d-2)\tau<\frac\beta2,
 \qquad HR\le Q,
 \qquad D\le N^\beta
\]
for large $N$; the first inequality follows directly from
$\tau=\beta\gamma/[1000k(d+1)]$.
By Lemma~\ref{transfer:specialization}, their Fourier coefficient
formulas for $q(\xi)\le H$ restrict the underlying integers to classes
$a'_1,a'_2\pmod D$.
The difference $a'_2-a'_1$ is admissible modulo $M_0$ by goodness.
It is admissible at each prime in $B$ because every conditioning of
positive probability is in the cycle support of the nonexceptional law.
It is therefore admissible modulo $D$.

The remaining frequencies of denominator at most $H$ are precisely the
set $\Xi$ in Proposition~\ref{kernel:estimate}.
Indeed, their denominators are squarefree and coprime to $D$.
Conversely, a squarefree denominator at most $H$ and coprime to $D$
has no prime below $P$, since all those primes divide $M_0$, and no
prime in $B$; every prime divisor is at most $H\le Q$.

Define functions on $[N]$ by
\[
 J_i(x)=\1_A(x)\1_{I_i}(x)
             \1_{\{x\equiv a'_i\pmod D\}},\qquad i=1,2.
\]
They are bounded by one and have ordered supports.
Every difference from the first support to the second is positive,
so a zero polynomial value cannot be such a difference;
\eqref{eq:24} excludes all nonzero values of $h_l(\N_+)$.
Thus all hypotheses of Proposition~\ref{kernel:estimate} hold at scale
$N$.  The Fourier normalizations in \eqref{eq:25} are
\[
 \widehat g_i(\xi)
 =\frac{DN}{|I_i|}\widehat{J_i}_{[N]}(\xi)
 \qquad(q(\xi)\le H).
\]
Since \eqref{eq:16} has a factor $D$ outside its absolute value,
the low-denominator part of \eqref{eq:30} is at most
\[
 C H^{-\gamma/2}D\frac{N^2}{|I_1||I_2|}
 \ll H^{-\gamma/2}D N^{2\tau}.
\]

Both parts of \eqref{eq:30} are consequently bounded, uniformly in the
conditioning and in the term of the character expansion, by
\[
 O\!\left(N^{-\beta\gamma/4+(d+2)\tau+o(1)}\right).
\]
Averaging the conditioning costs no factor, and summing the character
expansion costs at most $O(N^{2d\tau})$.
The final exponent satisfies
\[
 \frac{\beta\gamma}{4}-(3d+2)\tau-\sigma
 =\tau\left(250k(d+1)-(3d+2)-\frac\gamma4\right)>0.
\]
Thus the resulting error is $O(N^{-\sigma})$.
The same-interval contribution and the truncation error already satisfy
this bound.  Finally, the interval assignments were averaged with
probabilities summing to one, so the uniform estimate proves
\eqref{eq:29}.
\end{proof}
